\documentclass[a4paper,10pt]{article}
\usepackage[francais]{babel}
\usepackage{url}

\selectlanguage{francais}
\usepackage[utf8]{inputenc}
\usepackage[T1]{fontenc}
\usepackage{babel}
%\usepackage{graphicx}
\usepackage{xspace}
\usepackage{color}
\usepackage{setspace}
\usepackage{geometry}
%\usepackage{helvet}
\usepackage{calc}
\usepackage{ifthen}


\pagestyle{empty}

%%%%%
%%%%%

% Package pour la gestion des couleurs
%
\usepackage{color}

% Les commandes pour passer de la version nb à la version couleur 
% (et réciproquement)
%
\newboolean{enNoir}
\setboolean{enNoir}{true}
\newboolean{enCouleur}
\setboolean{enCouleur}{false}
\newcommand{\enNoir}{\definecolor{couleur}{rgb}{0,0,0}
  \setboolean{enCouleur}{false}
  \setboolean{enNoir}{true}}
\newcommand{\enGris}{\definecolor{couleur}{rgb}{0,0,0}
  \setboolean{enCouleur}{false}
  \setboolean{enNoir}{false}}
% Définition du bleu ULP
\newcommand{\enCouleur}{\definecolor{couleur}{rgb}{0,0.2,0.6}
  \setboolean{enCouleur}{true}
  \setboolean{enNoir}{false}}
% La version par défaut en en noir
%
\enNoir

%%%%%%
% Gestion de pdf
%%%%%%
%--------------------------------------
%pour compiler avec latex ou pdflatex
\newif\ifpdf
\ifx\pdfoutput\undefined
  \pdffalse
\else
  \pdfoutput=1
  \pdftrue
\fi
%--------------------------------------
%pour eviter de generer des fonts de type 3
\ifpdf 
\usepackage{pslatex}
\usepackage[pdftex]{graphicx}
%\DeclareGraphicsExtensions{.gif}
\else
\usepackage[dvips]{graphicx}
\fi
%--------------------------------------


\addtolength{\parskip}{\baselineskip}

\newlength{\firstcolumnwidth}
\setlength{\firstcolumnwidth}{1.65in}
\newlength{\inbetween}
\setlength{\inbetween}{5mm}
\newlength{\headermargin}
\setlength{\headermargin}{49mm}

\geometry{left=0.47in+\firstcolumnwidth+\inbetween,right=1in,top=15mm,bottom=20mm,noheadfoot}

\newenvironment{letter}[3][]
{
\thispagestyle{empty}

\newsavebox{\dimargin}
\sbox{\dimargin}{%
  \begin{minipage}[t]{\firstcolumnwidth}
    \vspace{0pt}
    \sffamily
    \begin{flushright}
      %%% Entête du DI
      \newsavebox{\dihead}
      \sbox{\dihead}{%
        \begin{minipage}{\linewidth}
          \flushright
          {\bfseries
            \color{couleur}
            \begin{spacing}{1}
              \large{Laboratoire}
              
              \large{des sciences de l'ingénieur, }
            
              \large{de l'informatique et} 

              \large{de l'imagerie (Icube)} 
        
              \normalsize{UMR 7537} 
            \end{spacing}
            }
          \vspace{-2mm}
          
          \rule{36.5mm}{0.5pt}
          
%          \vspace{-4mm}
          %%% Inclusion du Logo ULP
          \ifthenelse{\boolean{enNoir}}
          {
            \flushright{\includegraphics[height=19mm]{logo-uds-couleur-800x433}}
          }
          {
            \ifthenelse{\boolean{enCouleur}}
            {
              \includegraphics[height=19mm]{}
              }
            {
              \includegraphics[height=19mm]{logo-uds-couleur-800x433}
              }
            }
          %%% Fin inclusion logo
        \end{minipage}%
        }
      \newlength{\diheadheight}
      \newlength{\diheaddepth}
      \newlength{\diheadspace}
      \settoheight{\diheadheight}{\usebox{\dihead}}
      \settodepth{\diheaddepth}{\usebox{\dihead}}
      \setlength{\diheadspace}{\diheadheight+\diheaddepth}
      %
      %
      %%% Coordonnées de l'expéditeur
      \newsavebox{\sender}
      \sbox{\sender}{%
        \begin{minipage}{\linewidth}
          \flushright
          \begin{spacing}{1.5}
            
            \textbf{\small{\fromname}}

            \textbf{\footnotesize{\fonction}}
            
            \scriptsize{Tel. : \phoneespla}
            
            \scriptsize{\phoneillkirch}
            
            \scriptsize{\emailadd}
          \end{spacing}
        \end{minipage}%
        }
      \newlength{\senderheight}
      \newlength{\senderdepth}
      \newlength{\senderspace}
      \settoheight{\senderheight}{\usebox{\sender}}
      \settodepth{\senderdepth}{\usebox{\sender}}
      \setlength{\senderspace}{\senderheight+\senderdepth}
      %
      %
      %%% Coordonnées du DI
      \newsavebox{\difoot}
      \sbox{\difoot}{%
        \begin{minipage}{\linewidth}
          \flushright
          \begin{spacing}{1.25}
            \footnotesize
            
            \textbf{Laboratoire Icube}
            
            
            {Pôle API -Bd S. Brant}
            
            {F-67400 Illkirch}
            
            {Tél. : +33 (0) 3 90 24 45 42}
            
            {Fax : +33 (0) 3 90 24 45 47}
            
            {http\string://icube.unistra.fr}
          \end{spacing}
        \end{minipage}%
        }
      \newlength{\difootheight}
      \newlength{\difootdepth}
      \newlength{\difootspace}
      \settoheight{\difootheight}{\usebox{\difoot}}
      \settodepth{\difootdepth}{\usebox{\difoot}}
      \setlength{\difootspace}{\difootheight+\difootdepth}
      %
      %
      \newlength{\devantsender}
      \setlength{\devantsender}{57mm}
%
      \newlength{\devantdifoot}
%
      \setlength{\devantdifoot}{\textheight-\diheadspace-\devantsender-\senderspace-\difootspace-8mm-13pt}
%
      \vspace{8mm}

      \usebox{\dihead}

      \vspace{\devantsender}

      \usebox{\sender}

      \vspace{\devantdifoot}

      \usebox{\difoot}
    \end{flushright}
  \end{minipage}%
}

\begin{spacing}{1}

  \noindent
  \begin{picture}(0,0)
    \setlength{\unitlength}{\firstcolumnwidth+\inbetween}
    \put(-1,0){\usebox{\dimargin}}
  \end{picture}%
%
\begin{minipage}[t]{\headermargin}

\vspace{26mm}

~
\end{minipage}%
\begin{minipage}[t]{\linewidth-\headermargin}
  %Strasbourg, le 30 octobre 2012
  Strasbourg, le \today
\end{minipage}

\noindent
\begin{minipage}[t]{\headermargin}
\vspace{35mm}
~
\end{minipage}%
%-----------------------------------------
% Champ destinataire (#2=To/A , #3 <nom>)   
%-----------------------------------------
\ifthenelse{\equal{#3}{}}{%
\begin{minipage}[t]{\linewidth-\headermargin}

\vspace{2mm}

\end{minipage}
}%
{%
\begin{minipage}[t]{\linewidth-\headermargin} #2
  
  

  \vspace{2mm}

  #3
}
\end{minipage}
\ifthenelse{\equal{#1}{}} {} {Objet : {#1}}
\vspace{20mm}
}
{
%\vspace{20mm}

\newlength{\taillesignature}
\settowidth{\taillesignature}{\fromname}
\newlength{\possignature}
\setlength{\possignature}{0.667\linewidth-0.5\taillesignature}

\hspace{\possignature}\fromname


\end{spacing}
}



\newcommand*{\name}[1]{\def\fromname{#1}}
\name{}
\newcommand*{\fonction}[1]{\def\fonction{#1}}
\newcommand*{\telephoneesplanade}[1]{\def\phoneespla{#1}}
\telephoneesplanade{}
\newcommand*{\telephoneillkirch}[1]{\def\phoneillkirch{#1}}
\telephoneillkirch{}
\newcommand*{\email}[1]{\def\emailadd{#1}}
\email{}
\newcommand*{\objet}[1]{\def\emailadd{#1}}





\selectlanguage{francais}
\begin{document}


\name{St\'ephane \textsc{Genaud}}
\fonction{Professeur des Universités}
\telephoneillkirch{+33 (0)3688 54542}
\email{genaud@unistra.fr}

\begin{letter}[Recommendation de Marc Frincu pour une \emph{Starting Research Position}]%
{To}{Président d'INRIA}
 

J'ai rencontré pour  la première fois Marc Fr\^incu en  2007, alors qu'il venait
faire  internship INRIA  au centre  de  Nancy-Grand-Est dans  l'équipe projet  à
laquelle j'étais  rattaché. Sans  l'avoir encadré directement,  je sais  qu'il a
donné  entièrement  satisfaction  à  mes  collègues  qui  lui  ont  proposé  une
allocation \emph{CORDI}.   Marc a décliné cette  proposition et a fait  sa thèse
dans son  université d'origine, Timisoara,  en Roumanie.  J'ai retrouvé  Marc en
2012 : il  cherchait un postdoc, et à  la vue du travail qu'il  avait réalisé en
thèse, j'ai décidé de convertir un financement de thèse pour l'accueillir sur un
postdoc  à  Strasbourg (il  était  convenu  de commencer  par  12  mois puis  de
prolonger à 18 si cela convenait aux deux parties).

Lors  de son  postdoc, j'ai  pu  découvrir au  fur  et à  mesure l'autonomie  et
l'efficacité de Marc. Le financement du postdoc reposant sur un projet ANR, Marc
avait  dans  ce cadre  des  objectifs  précis.  Il  a  été  très rigoureux  dans
l'accomplissement de  ces objectifs, mais cela  ne l'a pas empêché  de voir plus
loin. Il  a mené de  front la finalisation  d'articles qu'il avait  débuté avant
d'arriver   à  Strasbourg,   étudié   la  possibilité   de   monter  un   projet
France-Roumanie  (ANR programme  blanc international),  initié des  travaux plus
théoriques sur des thèmes connexes à ceux  du projet ANR. Je lui avais également
fixé  l'objectif de  faire une  synthèse des  travaux relatifs  aux nôtres  pour
redéfinir  la direction  dans laquelle  on devait  orienter nos  recherches. Son
travail nous a très vite donné l'idée  d'en faire un survey du domaine, que nous
avons publié depuis.

 
Au final, son séjour (septembre  2012-septembre 2013) a été extrêment productif.
En termes  de rédaction d'articles,  Marc a  été le principal  contributeur d'un
papier  dans un  workshop de  la  conférence IPDPS,  de deux  articles de  revue
internationales  (notifications  d'acceptations  récentes), et  a  contribué  de
manière secondaire à un papier dans une conférence internationale très sélective
(CloudComm). Si son postdoc n'a duré qu'un an, ce n'est que parce qu'il a
anticipé sa recherche du postdoc suivant, et notre regret mutuel est de ne pas
avoir prolongé cette période au maxium.


Cette  activité intense  est dûe  à  l'autonomie de  Marc, qualité  que je  dois
souligner ici.   Quand Marc propose  une idée,  il a généralement  déjà beaucoup
travaillé avant de  la présenter, et la discussion est  ainsi très fructueuse et
déjà aboutie.  Marc a de plus  une grande qualité  d'écoute : il fait  tout pour
intégrer  à son  travail les  remarques ou  idées de  ses contradicteurs.   Marc
possède  également une  culture  bibliographique remarquable  dans les  domaines
qu'il a étudié  en recherche, acquise grâce à de  nombreuses lectures.  Alliée à
un esprit  de synthèse  et des aptitudes  rédactionnelles évidentes,  toutes ces
qualités font que la transcription des  travaux de recherche en articles se font
sans peine.

J'ai également évoqué  la proposition de monter un projet  ANR international que
Marc m'a fait 3  mois après son arrivée. Cette initiative  démontre pour moi ses
qualités de  leadership et d'esprit  d'iniative.  L'expérience acquise  dans son
postdoc à  USC (University of  Southern California) a probablement  renforcé ses
qualité de leadership, car il m'a  indiqué avoir un rôle d'encadrement, avec une
partie importante  sur la rédaction  de réponses  à appels d'offre  d'agences de
financement comme la NSF.

Pour toutes ces  raisons, je recommende au  plus haut point Marc  Frincu pour ce
poste \emph{Starting Research Position} à l'INRIA.


\end{letter}
\begin{flushright}
\includegraphics[width=.30\textwidth]{signgenaud.jpg}
\end{flushright}
%\newpage
%\bibliographystyle{plain}
%\bibliography{julien_blaise_ED}


\end{document}

\end{letter}
\begin{flushright}
\includegraphics[width=.20\textwidth]{signgenaud.jpg}
\end{flushright}
%\newpage
%\bibliographystyle{plain}
%\bibliography{julien_blaise_ED}


\end{document}
